% To allow compilation of the file
% !TeX root = ./../article.tex

\section{Introduction}

We study federated elections in which authorized local authorities conduct
secret voting and compute local tallies off-chain. Ballot secrecy is
already enforced locally and local committees are assumed to count correctly.
The weak point is the transition from local results to the global result: a
central coordinator can omit unfavorable tallies, misreport the aggregate, or
exploit early knowledge of partial results.

Publishing local tallies on a blockchain makes inclusion auditable but exposes
local outcomes, which can create retaliation, reputational pressure, or
strategic bargaining, especially in small electorates. Commit-and-reveal only
postpones that disclosure and still leaves a final-stage abort surface.

This motivates the central design choice of this paper: \emph{the local tally
should remain encrypted throughout the protocol}. Each authority publishes an
encrypted tally vector and a zero-knowledge proof that the hidden tally is
structurally valid and sums to its public electorate size. The contract accepts
one such submission per authority, fixes the transcript on-chain, computes the
homomorphic aggregate, and only the global final tally is decrypted.

Our design is closer to open-audit bulletin board systems such as
Helios~\cite{USENIX:Adida08} and to multi-authority election
protocols~\cite{EC:CraGenSch97} than to a ballot-level remote voting stack, but
adds three elements for private local tallies: validity proofs for encrypted
local results, an on-chain aggregate updated from accepted submissions, and a
proof-checked threshold opening transcript that reveals only the final
aggregate.

We instantiate encryption with curve-based additive ElGamal, which aligns the
homomorphic layer with Circom/Groth16 arithmetic. The work remains at the
local-tally level because the problem addressed here is final aggregation from
trusted local counts, not ballot-level remote voting.
